\chapter{Методология}
\label{chap:met}

\section{Требования и подход к разработке}

В данном разделе мы подробно описываем требования к функциональности языкового
сервера.

\subsection*{Требования}
\label{chap:met:requierements}

Языковой сервер должен поддерживать ключевые функции для разработки:
\begin{itemize}
  \item Отображение информации о символе при наведении (подсказка о типе)
  \item Переход к определению
  \item Поддержка синтаксической/семантической подсветки
  \item Диагностика исходного кода (парсинг, проверка типов)
\end{itemize}

Также важно предоставить не только исполняемый файл сервера, но и чёткие
инструкции по его интеграции с редактором.

\subsection*{Архитектура}

Основываясь на архитектуре серверов для OCaml и Semgrep (см.
\ref{sec:background:implementations}), была выбрана расширяемая и асинхронная
архитектура. Сервер реализован на OCaml для лёгкой интеграции с компилятором
Jasmin, который написан на этом же языке.

Реализация разделена на два компонента: сервер JSON-RPC и группа обработчиков
LSP. Сервер отвечает за низкоуровневые операции протокола, а обработчики
реализуют логику.

